\documentclass[a4paper,12pt]{article} % Tipo de documento y tamaño de letra.
\usepackage[utf8]{inputenc} % Para que reconozca caracteres especiales como tildes.
\usepackage[spanish, es-nodecimaldot]{babel} % Para definir el idioma y el formato de fechas en español.
\usepackage{amsmath, amssymb} % Paquetes para símbolos y ecuaciones avanzadas (Para hacer una ecuación con un recuadro).
\usepackage{ragged2e} % Para justificar textos.
\usepackage{multicol} % Para crear columnas.
\usepackage{slashed} % Para generar números o variables canceladas (Una línea diagonal sobre el número o variable).
\usepackage{graphicx} % Paquete para poner graficas.
\usepackage{float} % Para poder dejar la imagen en un punto fijo del documento.
\usepackage{longtable} % Para tablas largas.
\usepackage{cancel} % Para hacer cancelaciones en medio de operaciones.
\usepackage{xcolor} % Para el color y esas cosas.
\usepackage{apacite} % Paquete para citas APA.
\usepackage{fancyhdr} % Para personalizar encabezados y pies de página.
\usepackage{geometry} % Para ajustar márgenes.
\usepackage{setspace} % Para ajustar el interlineado.
\usepackage{natbib} % Para manejar las citas y que no queden pegadas y feas :P.
\usepackage{hyperref} % Para que los enlaces funcionen con su URL.
\usepackage[version=4]{mhchem}  % Versión recomendada para mayor funcionalidad.
\usepackage{chemfig} %También para hacer estructuras químicas.
\usepackage{array} % Para mejor control de las columnas.
\usepackage{siunitx} % Para unidades físicas.
\usepackage{tabularx}
\usepackage[table]{xcolor}
\usepackage{booktabs}

\pagestyle{fancy}
\fancyhf{} % Limpia encabezados y pies de página predeterminados
\fancyhead[L]{\textbf{CIENTÍFICOS DE LA HISTORIA.}} % Encabezado izquierdo
\fancyhead[R]{\textbf{HISTORIA Y DIDÁCTICA.}} % Encabezado derecho
\fancyfoot[C]{\thepage} % Número de página en el centro del pie

\geometry{headheight=12.96115pt, top=1in, bottom=1in, left=1in, right=1in}
\begin{document}


\begin{figure}[H]
	\centering
	\includegraphics[width=0.5\textwidth]{Unir.jpg}
\end{figure}

\vspace{\baselineskip}


\begin{center}
	\textbf{\Large } \\[0.3cm]
	\vspace{\baselineskip}
	\textbf{} \\[0.4cm]
	\textbf{Grado en Física.} \\[0.8cm]
	\textbf{Enric Navarro Cuadro.} \\[0.4cm]
	\textbf{Liz ángel Núñez Torres.} \\[1 cm]
	\textbf{Universidad Internacional de La Rioja.} \\[1 cm]
	\textbf{22 de Marzo 2026.}
\end{center}

\newpage
\tableofcontents


\newgeometry{top=2cm, bottom=2cm, headheight=0cm, footskip=0cm}

\arrayrulecolor{blue} % Color de los bordes de la tabla
\setlength{\arrayrulewidth}{1.5pt} % Grosor de los bordes



\section{Grandes científicos de la geometría y astronomía.}

\vspace{\baselineskip}

\vspace{\baselineskip}

\vspace{\baselineskip}

\begin{table}[H]
    \centering
    \begin{tabular}{ll}
        \toprule
        \textbf{ASTRÓNOMOS.} & \textbf{GEÓMETRAS.} \\
        \midrule
        Eudoxo de Cnido. & Tales de Mileto. \\
        Aristarco de Samos. & Pitágoras. \\
        Hiparco de Nicea. & Euclides. \\
        Claudio Ptolomeo. & Eratóstenes. \\
		Nicolás Copérnico. & Apolonio de Perga. \\
        \bottomrule
    \end{tabular}
    
    \label{tab:elegante}
\end{table}

\vspace{\baselineskip}

\vspace{\baselineskip}

\vspace{\baselineskip}

\vspace{\baselineskip}

\vspace{\baselineskip}

\section{Astrónomos.}
\vspace{\baselineskip}
\vspace{\baselineskip}


\subsection{Eudoxo de Cnido (c. 408 a. C. – c. 355 a. C.) }
\begin{justify}
	Es considerado el fundador de la astronomía matemática, pues fue el primero en desarrollar un modelo geométrico para explicar los
	movimientos aparentes de los planetas, el Sol y la Luna mediante un sistema de esferas homocéntricas. Su obra sentó las bases sobre las que construirían Calipo, Aristóteles y, posteriormente, toda la tradición astronómica griega.
\end{justify}

\begin{justify}
	Propuso que cada astro se movía mediante una combinación de esferas concéntricas con centro en la Tierra, cada una con su propio eje de rotación inclinado y su propia velocidad. Para el Sol y la Luna utilizó tres esferas cada uno; para los planetas 
	(Mercurio, Venus, Marte, Júpiter y Saturno) utilizó cuatro esferas cada uno, sumando un total de 27 esferas. Este modelo lograba explicar cualitativamente fenómenos como el movimiento retrógrado de los planetas y la latitud variable de la Luna.
\end{justify}

\begin{justify}
	Demos
	 principio de explicar movimientos complejos mediante la composición de movimientos circulares simples se convirtió en el paradigma central de la astronomía durante más de dos milenios.
\end{justify}

\newpage

\subsection{ Aristarco de Samos (c. 310 a.C. -- c. 230 a.C.)}
\begin{justify}
	En su única obra conservada, \textbf{``Sobre los tamaños y las distancias del Sol y la Luna''}, Aristarco desarrolló un método 
	geométrico para calcular las distancias relativas de estos dos astros. Basándose en observaciones de los cuartos de Luna 
	(cuando la Luna aparece exactamente medio iluminada), determinó que el ángulo Sol-Tierra-Luna es de 87° (en realidad es 89°51'). 
\end{justify}

\begin{justify}
	A partir de este dato, mediante trigonometría elemental, concluyó que la distancia al Sol es entre 18 y 20 veces mayor que la distancia a la Luna 
	(el valor real es unas 400 veces). Aunque sus mediciones angulares eran imprecisas debido a la dificultad de determinar exactamente el momento del cuarto de Luna, 
	su método geométrico era impecable y sentó las bases para la astronomía de distancias.
\end{justify}

\begin{justify}
	Aristarco fue el primero en postular que la Tierra gira alrededor del Sol y no al revés.
\end{justify}

\subsection{ Hiparco de Nicea (c. 190 a.C. -- c. 120 a.C.)}
\begin{justify}
	Para sus cálculos astronómicos, desarrolló la primera tabla de cuerdas (antecesora de las funciones seno), 
	que tabulaba las longitudes de las cuerdas correspondientes a ángulos centrales en un círculo. Dividió la circunferencia en 360 grados
	(heredando este sistema de los babilonios) y desarrolló métodos para calcular cuerdas de ángulos suma, diferencia y mitad, estableciendo 
	las bases de la trigonometría plana y esférica que serían esenciales para la astronomía durante los siguientes 1.500 años.
\end{justify}

\begin{justify}
	\begin{itemize}
		\item Descubrió que las posiciones de las estrellas no eran fijas, sino que experimentaban un desplazamiento gradual hacia el este a lo largo de los siglos. 
		Determinó que este fenómeno, llamado precesión de los equinoccios, tenía una tasa de aproximadamente 1° por siglo (el valor real es 1° cada 72 años). 
		
		\item Compiló el primer catálogo estelar sistemático de la historia occidental, registrando las posiciones y magnitudes (brillo) de unas 850 a 1.080 estrellas. 
		Para clasificarlas, inventó un sistema de seis magnitudes aparentes (del 1 al 6) que, con modificaciones, sigue utilizándose en la astronomía moderna.

		\item Refinó las mediciones del año solar, determinando su duración con una precisión asombrosa para la época: 365 días, 5 horas, 55 minutos y 12 segundos (solo unos 6 minutos más largo que el valor real).
		
		\item Utilizando observaciones de eclipses solares totales desde diferentes localizaciones, Hiparco calculó la distancia a la Luna con una precisión sin precedentes, estimándola en 59 radios terrestres (el valor real es aproximadamente 60).
	\end{itemize}
\end{justify}


\newpage 

\subsection{ Claudio Ptolomeo (c. 100 d.C. -- c. 170 d.C.)}
\begin{justify}
	Su obra \textbf{Almagesto}, escrita hacia 150
	d. C., fue el manual astronómico de referencia en el mundo islámico y en Europa hasta comienzos
	del siglo XVII.
\end{justify}
\begin{justify}
	Formuló el sistema geocéntrico ptolomaico: que situaba la Tierra inmóvil en el centro, y el Sol, la
	Luna y los planetas moviéndose mediante combinaciones de movimientos circulares.
	Para explicar irregularidades aparentes, como la retrogradación de los planetas, usó deferentes,
	epiciclos y el ecuante. Consiguió un modelo matemático bastante preciso para predecir posiciones
	planetarias.
\end{justify}
\begin{justify}
	Aunque hoy sabemos que el sistema geocéntrico era falso, su teoría funcionaba suficientemente bien
	como herramienta de cálculo y por eso dominó la astronomía durante unos 1.400 años.
	También incluyó tablas trigonométricas y un catálogo de unas 1.000 estrellas, además de
	procedimientos para calcular posiciones astronómicas en fechas concretas.
\end{justify}

\subsection{ Nicolás Copérnico (1473 -- 1543)}

\begin{justify}
	Su obra \textbf{De revolutionibus orbium coelestium} (Sobre las revoluciones de las esferas celestes), 
	publicada en el año de su muerte, marcó el inicio de la revolución científica al restaurar y 
	demostrar matemáticamente el modelo heliocéntrico propuesto por Aristarco de Samos dieciocho siglos antes. 
	Su sistema desplazó a la Tierra del centro del universo y la colocó como un planeta más que orbita alrededor del Sol, 
	desafiando 1.400 años de tradición aristotélico-ptolemaica.
\end{justify}

\begin{justify}
	\begin{itemize}
		\item Estableció definitivamente el orden de los planetas en función de sus períodos orbitales, algo que en el sistema geocéntrico era arbitrario o basado en criterios no estrictamente astronómicos. 
		\item Su modelo permitió calcular las distancias relativas de los planetas al Sol con una precisión muy superior a la del sistema ptolemaico. Al situar al Sol en el centro y a la Tierra como un planeta más, 
		pudo utilizar observaciones de elongaciones máximas de los planetas interiores y de oposiciones de los exteriores para calcular las razones de sus distancias a la del Sol-Tierra, estableciendo una escala coherente para todo el sistema solar.
		\item Introdujo un nuevo modelo que explicaba con mayor precisión las irregularidades del movimiento lunar (como la variación y la evección), basándose en observaciones propias y en datos de astrónomos árabes. 
	\end{itemize}
\end{justify}

\newpage

\section{Geómetras.}
\vspace{\baselineskip}
\vspace{\baselineskip}


\subsection{Tales de Mileto (c. 624 a.C. -- c. 546 a.C.)}
\begin{justify}
	Es considerado una figura
	fundamental de geometría griega, y se le atribuyen varios teoremas elementales y aplicaciones
	prácticas de semejanza y medida, no obstante sigue siendo difícil diferenciar la leyenda, del hombre
	y lo que realmente demostró.
	Entre sus proposiciones están:
\end{justify}

\begin{justify}
	\begin{itemize}
		\item El diámetro divide el círculo en dos partes iguales. Si se traza una línea que atraviese el centro, se
		crean dos semicircumferencias de igual tamaño, por lo tanto el centro está a la misma distancia de
		todos los puntos de la circumferencia.

		\item Los ángulos de la base de un triángulo isósceles son iguales. Asumiendo la idea general de que en
		un triangular, a lados iguales, ángulos opuestos iguales.

		\item  Si dos rectas se cortan, los ángulos opuestos por el vértice son iguales. Cuando dos rectas se
		cruzan formando una \(\mathrm{X}\), los ángulos enfrentados son iguales dos a dos.

		\item Un triángulo queda determinado por su base y los dos ángulos adyacentes a esa base. Si se conoce
		un lado y los ángulos de sus extremos, el triángulo queda fijado por estos.

		\item El ángulo inscrito en una semicircunferencia es recto. La atribución más famosa, dice que si se
		toma un diámetro \(\mathrm{AB}\) de una circumferencia, y se toma cualquier punto \(\mathrm{C}\) dentro de la
		semicircumferencia, el ángulo \(<\mathrm{ACB}\) mide \(\SI{90}{\degreeCelsius}\).
	\end{itemize}
\end{justify}

\subsection{Pitágoras (c. 569 a.C. -- c. 475 a.C.)}
\begin{justify}
	Aunque el enunciado era conocido en civilizaciones anteriores (babilonios, egipcios, chinos), los pitagóricos fueron los 
	primeros en demostrarlo de manera rigurosa y general: \textbf{``En todo triángulo rectángulo, el cuadrado construido sobre la hipotenusa 
	es igual a la suma de los cuadrados construidos sobre los catetos"}. Esta demostración, basada en consideraciones de áreas,
	convirtió el teorema en el resultado geométrico más influyente de la historia, con aplicaciones que se extienden a todas las 
	ramas de las matemáticas y la física.
\end{justify}
\begin{justify}
	La aplicación del \textbf{teorema de Pitágoras} al triángulo rectángulo isósceles (de catetos iguales a 1) condujo al descubrimiento de que 
	la diagonal y el lado son magnitudes inconmensurables: su razón no puede expresarse como cociente de números enteros. 
\end{justify}
\begin{justify}
	Este hallazgo, atribuido a \textbf{Hipaso de Metaponto}, demostró que la relación geométrica entre segmentos podía escapar a la aritmética
	de los números enteros, lo que supuso una crisis epistemológica para la doctrina pitagórica ("todo es número"), pero también abrió 
	el camino hacia el concepto de número irracional y la geometría como fundamento de las matemáticas.
\end{justify}

\subsection{Euclides (c. 325 a.C. -- c. 265 a.C.)}
\begin{justify}
	Su obra más famosa es \textbf{Los Elementos}, un tratado
	matemático, sobre todo de geometría, donde reúne definiciones, axiomas, postulados y teoremas
	demostrados paso a paso. Fue durante siglos el libro básico para estudiar matemáticas. Convirtió la
	geometría en un sistema deductivo riguroso, con organización global, nivel de demostración y
	diseño de conjunto.
\end{justify}
\begin{justify}
	En sus libros, Euclides dio forma rigurosa a la teoría de la proporción y la semejanza, así como
	también sistematizó toda la geometría plana elemental.
	También desarrolló lo que hoy llamamos «álgebra geométrica», y expresó construcciones como la
	sección áurea, y un equivalente a la ley de cosenos.
\end{justify}
\begin{justify}
	En su quinto postulado, o de las paralelas, formuló el marco de la geometría euclidiana propiamente
	dicha, y el cuestionamiento e intento de demostración de este durante siglos, dio lugar a separarlo de
	la geometría no euclidiana (elípticas e hiperbólicas).
\end{justify}

\subsection{Eratóstenes (c. 276 a. C. – c. 194 a. C.)}
\begin{justify}
	En su obra \textbf{Geographica} (hoy perdida, pero conocida por resúmenes posteriores), 
	Eratóstenes aplicó principios geométricos a la descripción de la Tierra. Fue el primero en trazar un sistema de paralelos y meridianos 
	(líneas de latitud y longitud) para construir mapas con bases matemáticas, dividiendo el mundo conocido en secciones climáticas.
\end{justify} 

\begin{justify}
	Calculó la oblicuidad de la eclíptica 
	(la inclinación del eje terrestre) en 11/83 de circunferencia (unos 23°51'), un valor notablemente preciso para la época. También estimó distancias astronómicas, como 
	la distancia entre la Tierra y el Sol, utilizando métodos geométricos basados en observaciones de eclipses lunares.
\end{justify}

\begin{justify}
	Calculó la distancia de la Tierra a la Luna y al Sol, así como los diámetros de estos cuerpos, utilizando observaciones de eclipses y principios de geometría de triángulos semejantes.
	Aunque sus cifras no fueron tan precisas como la de la circunferencia terrestre, su método geométrico sentó las bases de la astronomía matemática.
\end{justify}

\subsection{Apolonio de Perge (c. 262 a.C. -- c. 190 a.C.)}
\begin{justify}
	Su obra \textbf{Almagesto}, escrita hacia 150
	d. C., fue el manual astronómico de referencia en el mundo islámico y en Europa hasta comienzos
	del siglo XVII.
\end{justify}
\begin{justify}
	Formuló el sistema geocéntrico ptolomaico: que situaba la Tierra inmóvil en el centro, y el Sol, la
	Luna y los planetas moviéndose mediante combinaciones de movimientos circulares.
	Para explicar irregularidades aparentes, como la retrogradación de los planetas, usó deferentes,
	epiciclos y el ecuante. Consiguió un modelo matemático bastante preciso para predecir posiciones
	planetarias.
\begin{justify}
\end{justify}
	Aunque hoy sabemos que el sistema geocéntrico era falso, su teoría funcionaba suficientemente bien
	como herramienta de cálculo y por eso dominó la astronomía durante unos 1.400 años.
	También incluyó tablas trigonométricas y un catálogo de unas 1.000 estrellas, además de
	procedimientos para calcular posiciones astronómicas en fechas concretas.
\end{justify}


\newpage

\section{Bibliografía.}
\setlength{\bibsep}{12pt} % Ajusta la separación entre referencias
\bibliographystyle{apacite} % Estilo APA
\bibliography{citas} % Archivo de referencias
\nocite{*}


\newpage

\section{Revisión.}

\begin{figure}[H]
	\centering
	\includegraphics[width=1.1\textwidth]{revision.png}
\end{figure}

\end{document}